
\begin{frame}[fragile]{Séance du 3 juin 2025}

Nous allons étudier de manière complète les phases
d'une mise en œuvre d'Elastic Search. Installation, chargement de données,
interrogation et passage à l'échelle.

\vfill

Pour l'installation, deux possibilités\,:

\vfill 

\begin{redbullet}
  \item Créer un compte d'essai sur \url{https://cloud.elastic.co/}. Facile, mais 
     pas de possibilité de tester les fonctions multi-serveur
  \item Installer localement avec Docker\,: très utile et instructif
\end{redbullet}

\vfill

Dans ce qui suit\,: description de la seconde solution (recommandée).

\end{frame}

\begin{frame}[fragile]{Lançons-nous avec Docker}

En bref\,: Docker nous permet d'installer localement, avec
très peu de ressources, des serveurs virtuels (ou \alert{conteneurs}).


\vfill 
En créant plusieurs serveurs du même type (par exemple ElasticSearch) et
en les faisant communiquer on obtient un \alert{système distribué}.

\vfill 
Pour une introduction détaillée, voir \url{http://b3d.bdpedia.fr/docker.html}

\end{frame}



\begin{frame}{Système distribué: rappels}

\begin{redbullet}
  \item  un \alert{serveur machine} est un ordinateur, tournant sous un système d'exploitation,
    et connecté en permanence au réseau via des ports; il a une IP (adresse internet)
  \item   un \alert{serveur logiciel} est un processus exécuté en tâche de fond d'un serveur machine qui
    communique avec des \alert{clients (logiciels)} via un port particulier.
  \item Une \alert{machine virtuelle} est un programme qui simule, sur une machine hôte,
     un autre ordinateur.
   \item Un \alert{client} (logiciel) est un programme qui communique avec un serveur (logiciel)
\end{redbullet}
    
\begin{defblock}{Exemple}
Un serveur web est un processus (Apache par exemple) qui communique sur le
port 80 d'un serveur machine. Si ce serveur machine a pour IP 163.12.9.10, alors tout
client web peut s'adresser au serveur web à l'adresse 163.12.9.10:80.
\end{defblock}

\end{frame}


\begin{frame}{Illustration}

Voici un exemple de système distribué géré par Docker.

\vfill

\figSlideWithSize{Serveurs}{12cm}

\end{frame}

\begin{frame}{Docker}

Docker permet de créer des \alert{conteneurs}, des machines virtuelles extrêmement légères.
\vfill

\figSlideWithSize{schema-docker}{9cm}

\vfill

Existe sous Linux, ainsi que MacOS et Windows.
\end{frame}


\begin{frame}{Les images}

Image = système pré-configuré, prêt à l'emploi pour un service particulier (par exemple un serveur
MySQL, ou MongoDB).

\vfill

\figSlideWithSize{images-docker}{8cm}

\vfill
Une image s'installe très facilement. C'est une sorte de moule pour créer
(``instancier'')  des conteneurs.

\end{frame}


\begin{frame}{Comment communiquer avec un conteneur ?}

J'ai lancé un conteneur ElasticSearch. Mon serveur tourne, sur le port 
9200, \alert{dans le conteneur}.

\vfill

\alert{Question:} comment peut-on communiquer avec ce serveur?

\vfill

\alert{Réponse}: \emph{il faut associer un port de la machine Docker avec le port 
9200 du conteneur} (\textit{port forwarding}).

\vfill

Par exemple, 
\begin{redbullet}
 \item je vais associer  le port 9200 de la machine Docker avec le port 9200 du conteneur \texttt{elastic1}
 \item je vais associer le port 9201 de la machine Docker avec le port 9200 du conteneur \texttt{elastic2}
\end{redbullet}


\end{frame}

\begin{frame}{Installons Docker}

Rendez-vous sur \url{https://docker.com}

\vfill

\begin{redbullet}
	\item Créez un compte 
	\item Récupérez l'application \textit{Docker desktop}. Lancez-la
\end{redbullet}

\vfill

\figSlideWithSize{docker-dashboard.jpg}{8cm}

\end{frame}


\begin{frame}[fragile]{Lançons un serveur (conteneur) ElasticSearch}


Ouvrez un terminal et lancez la commande suivante:

\begin{minted}{bash}
docker run -d --name es1 -p 9200:9200  elasticsearch:8.13.0
\end{minted}

\vfil
\begin{redbullet}
	\item On demande à instancier la version 8.13 d'Elastic Search  
	\item Le nom du conteneur est \texttt{es1}
	\item Le port 9200 du conteneur est associé au port 9200 de l'hôte
	\item Le processus s'exécute en détaché 
\end{redbullet}

\end{frame}


\begin{frame}[fragile]{Vérifier que le conteneur est bien actif}

Lancez le \textit{Docker desktop} et consultez la liste des conteneurs.

\vfill

\figSlideWithSize{docker-avec-es1}{11cm}
\end{frame}



\begin{frame}[fragile]{Il reste à créer un compte}

On crée un compte \texttt{elastic} avec la commande suivante:


\begin{minted}{bash}
docker exec -it es1 /usr/share/elasticsearch/bin/elasticsearch-reset-password -u elastic
\end{minted}

\vfill

Il nes reste plus qu'à se connecter...
\end{frame}

\begin{frame}{Il nous faut un client}

Pour communiquer avec le serveur, il nous faut une application
cliente. Nous allons utiliser ElasticVue (\url{https://elasticvue.com/})

\vfill

Installez l'application et connectez-vous au serveur sur le port 9200 
(\url{https://localhost:9200})
\vfill

\figSlideWithSize{elasticvue.png}{11cm}

\end{frame}




\begin{frame}{La création d'un index}


\figSlideWithSize{es-create-index.png}{8cm}
\end{frame}


